\scanpage{175}

The product of the last three terms on the right-hand side is not encouraging:
\[
\underline{u}'f\underline{u}'^{-1}
=\mu'\deleted{\unsure{\frac{1}{\mu'}}}
\begin{pmatrix}
\varepsilon'ru'^2+st'u'+tt'u'+\varepsilon'ut'^2
&-\varepsilon'rt'u'+su'^2-tt'^2+\varepsilon'uu't'\\
-\varepsilon'ru't'\unsure{-}st'^2+tu'^2+\varepsilon'uu't'
&-\varepsilon'rt'^2\unsure{-}su't'-tu't'+\varepsilon'uu'^2
\end{pmatrix},
\]
but the product with $f'$ simplifies miraculously, thanks to the relation
\[
\unsure{t'^2}+u'^2=\varepsilon'/\mu'
\]
and we find
\[
\deleted{\unsure{(17)}}\qquad
f''=\begin{pmatrix}r''&s''\\t''&u''\end{pmatrix}
=\deleted{\unsure{\frac{1}{\mu'}}}
\begin{pmatrix}
\unsure{\varepsilon|*}\ \text{mess}&\text{mess}\\
\varepsilon'tu'+ut'&-\varepsilon'tt'+uu'
\end{pmatrix}
\]
whence
\[
\tag{18}
t''=\frac{1}{\mu'}(\varepsilon'tu'+ut')
\qquad
u''=\frac{1}{\mu'}(-tt'+uu')
\]
and hence (taking account of \inserted{$\mu''=\mu\mu'$,} $\varepsilon''=\varepsilon\varepsilon'$)
\[
\deleted{\unsure{(19)}}\qquad
\underline{u}''=\mu\mu'\deleted{\unsure{\mu'}}
\begin{pmatrix}
-\varepsilon'tt'+uu'&\varepsilon tu'+\varepsilon\varepsilon'ut'\\
-\varepsilon'tu'-ut'&-\varepsilon\varepsilon'tt'+\varepsilon\varepsilon'uu'
\end{pmatrix}
\]
and the formula to be checked \unsure{$u$},
\[
\underline{u}''=\underline{u}'\underline{u}
\]
reduces to
\[
\mu\mu'\begin{pmatrix}
-tt'+uu'&\varepsilon tu'+\varepsilon\varepsilon'ut'\\
-\varepsilon'tu'\deleted{\unsure{\varepsilon'}}-ut'&-\varepsilon\deleted{\unsure{\varepsilon'}}tt'+\varepsilon\varepsilon'uu'
\end{pmatrix}
=\mu'\begin{pmatrix}u'&\varepsilon't'\\-t'&\varepsilon'u'\end{pmatrix}
\mu\begin{pmatrix}u&\varepsilon t\\-t&\varepsilon u\end{pmatrix}
\]
which is indeed true, OK !.

Here again, one feels that these calculations ought to fit into a carpet of linear algebra that I am too lazy to unroll in extenso. One would have to paraphrase the group \inserted{\unsure{Or rather}} $\widetilde{M}_{0,3}$ -- described by these systems $(g,h,f)$ -- in terms \inserted{$(\alpha,\beta,f)$} of the reductions of $g,h,f$

\scanpage{176}
\mspage{387}

in $\mathrm{Gl}(2,\widehat{\mathbb Z})^{\pm1\,\prime}_{\unsure{\mathfrak g}}$, then replacing $\widehat{\mathbb Z}$ by an arbitrary ring, and of course $\mathrm{Gl}(2,\widehat{\mathbb Z})'$ by $\mathrm{Gl}(2,\unsure{A|\widehat{\mathbb Z}})$ itself. Perhaps I will come back to this eventually....

Having thus constructed
\[
\varphi:\widetilde{M}_{0,3}\longrightarrow\mathrm{Gl}(2,\widehat{\mathbb Z})'=\mathrm{Gl}(2,\widehat{\mathbb Z})/\mathord{\pm1}
\]
we deduce from it, by inverse image of the \inserted{central} extension $\mathrm{Gl}(2,\widehat{\mathbb Z})$ of $\mathrm{Gl}(2,\widehat{\mathbb Z})'$ by \deleted{$\unsure{\mathbb Z}$} $\{\pm1\}$, a central extension of $\widetilde{M}_{0,3}$ by $\{\pm1\}$, which I denote by $\unsure{\widetilde{N}_{1,1}}$, and we have a homomorphism of extensions
\[
\tag{19}
\diagram{\begin{tikzcd}[ampersand replacement=\&]
1\arrow[r]\&\{\pm1\}\arrow[r]\arrow[d,equal]\&\widetilde{N}_{1,1}\arrow[r]\arrow[d,"\varphi"]\&\deleted{\unsure{\mathcal T}}\widetilde{M}_{0,3}\arrow[r]\arrow[d,"\varphi'"]\&1\\
1\arrow[r]\&\{\pm1\}\arrow[r]\&\mathrm{Gl}(2,\widehat{\mathbb Z})\arrow[r]\&\mathrm{Gl}(2,\widehat{\mathbb Z})'\arrow[r]\&1
\end{tikzcd}}
\]

\deleted{On the other hand, denoting by \deleted{$\unsure{\widetilde{N}_{1,1}}$} $\unsure{\widetilde{\mathcal T}_{1,1}^{+}|\widehat{\mathcal T}_{1,1}^{+}}$ the inverse image of \deleted{$\unsure{\widetilde{N}_{1,1}}$} $\widehat{\mathcal G}_{0,3}^{+}\simeq\unsure{\widehat{\mathrm{Sl}(2,\mathbb Z)}'}\subset\widetilde{M}_{0,3}$ in the extension \deleted{$\unsure{\mathcal T}$} $\widetilde{N}_{1,1}$ of $\widetilde{M}_{0,3}$ by $\pm1$, we will have
\[
\diagram{\begin{tikzcd}[ampersand replacement=\&]
1\arrow[r]\&\unsure{\widehat{\mathcal T}_{1,1}^{+}|\widetilde{\mathcal T}_{1,1}^{+}}\arrow[r]\&\widetilde{N}_{1,1}
\end{tikzcd}}
\]
}

\scanpage{177}

One should check that the homomorphism $\varphi'$ we have just constructed coincides on $M_{0,3}\subset\widetilde{M}_{0,3}$ with the well-known ``modular'' homomorphism of \S\ 39 \unsure{(1)} -- it suffices to check it on the subgroup $\Gamma_{\mathbb Q,\sigma}$ of $\widetilde{\Gamma}_{\sigma}$ corresponding to $\Gamma_{\mathbb Q}$ -- this should not be very difficult, and I will \unsure{assume it} here. This \deleted{\unsure{being}} assumed, we find that the inverse image of $M_{0,3}$ in \deleted{$\unsure{\mathcal T}$} $\widetilde{N}_{1,1}$ is none other (up to canonical isomorphism) than $N_{1,1}$ itself -- in any case the inverse image of $\widehat{\mathcal G}_{0,3}^{+}$ \deleted{\unsure{on}} is none other than $\widehat{\mathcal T}_{1,1}^{+}$ \inserted{$\simeq\mathrm{Sl}(2,\mathbb Z)^\wedge$}. We will then have (with this identification)
\[
\tag{20}
\diagram{\begin{tikzcd}[ampersand replacement=\&]
1\arrow[r]\&\begin{gathered}\widehat{\mathcal T}_{1,1}^{+}\\\unsure{\Vert}\ \mathrm{Sl}(2,\mathbb Z)^\wedge\end{gathered}\arrow[r]\arrow[d]\&\deleted{\unsure{\mathcal G}}\widetilde{N}_{1,1}\arrow[r]\arrow[d,"\varphi"]\&\widetilde{\Gamma}\arrow[r]\arrow[d,"\unsure{\chi}"]\&1\\
1\arrow[r]\&\mathrm{Sl}(2,\widehat{\mathbb Z})\arrow[r]\&\mathrm{Gl}(2,\widehat{\mathbb Z})\arrow[r,"\det"]\&\widehat{\mathbb Z}^{*}\arrow[r]\&1
\end{tikzcd}}
\]

\scanpage{178}
\mspage{388}

so we have \unsure{reconstructed|reconstructed} in terms of loop \unsure{automorphisms} of $\widehat{\Pi}_{0,3}$ the diagram (1) of \S\ 39, or rather a diagram \unsure{a} larger, into which that of \S\ 39 embeds -- it would still remain to characterize $\Gamma_{\mathbb Q}$ as a subgroup of $\widetilde{\Gamma}$.

The diagram deduced directly from the preceding theorem
\[
\tag{21}
\diagram{\begin{tikzcd}[ampersand replacement=\&]
1\arrow[r]\&\widehat{\mathcal G}_{0,3}^{+}\arrow[r]\arrow[d]\&\widetilde{M}_{0,3}\arrow[r]\arrow[d,"\varphi'"]\&\widetilde{\Gamma}\arrow[r]\arrow[d,"\chi"]\&1\\
1\arrow[r]\&\mathrm{Sl}(2,\widehat{\mathbb Z})'\arrow[r]\&\unsure{\mathrm{Gl}(2,\widehat{\mathbb Z})'}\arrow[r,"\det"']\&\widehat{\mathbb Z}^{\unsure{*}}\arrow[r]\&1
\end{tikzcd}}
\]
can be regarded as being deduced from the preceding one by reduction modulo the \deleted{centers} \inserted{central subgroups} $\{\pm1\}$ in the two rows.

What becomes in (20) of the splitting (3) of $\widetilde{M}_{0,3}$? Is there a privileged way, for \unsure{every} element of $\widetilde{N}_{1,1}$ -- i.e. a pair of a $u=(g,h,f)\in\widetilde{M}_{0,3}$, and a $\underline{u}_0\in\mathrm{Gl}(2,\widehat{\mathbb Z})$

\scanpage{179}

which lifts the element $\underline{u}=\varphi'_{\unsure{\sigma}}(u)$ of $\mathrm{Gl}(2,\widehat{\mathbb Z})'$, to ``normalize'' it modulo \deleted{\illegible{1}} composition with an element of $\mathrm{Sl}(2,\mathbb Z)^\wedge$, so that it has \deleted{a} (say) a commutation property with $\sigma_\infty$, of the type $u(\sigma_\infty)=\deleted{\varepsilon}\sigma_\infty$, \deleted{\illegible{1}} $\unsure{u_0|u}(\sigma_\infty)=\pm\sigma_\infty$?

Another way of posing the question, in terms of the \unsure{subgroup} $\widetilde{\Gamma}_\sigma\subset\widetilde{M}_{0,3}$, is the following: is there a canonical lifting $\widetilde{\Gamma}_\sigma\longrightarrow\widetilde{N}_{1,1}$, thus assigning to every $\underline{u}=(g,f)\in\widetilde{\Gamma}_\sigma$ (\unsure{That is} a $\underline{u}=\mu(t\sigma_\infty+u)\tau_\infty^\alpha\in\mathrm{Gl}(2,\widehat{\mathbb Z})'$\unsure{,|\}} \deleted{\unsure{modulo}} i.e. a $\underline{u}$ defined modulo sign) a representative $\underline{u}_0\in\mathrm{Gl}(2,\widehat{\mathbb Z})$ -- i.e. a way of removing this ambiguity of sign? In fact, if we replace $\widetilde{\Gamma}_\sigma$ by $\Gamma_{\mathbb Q,\sigma}$, the Standard \deleted{\illegible{1}}\inserted{\unsure{yogas}} \unsure{suggest} that the liftings in

\scanpage{180}
\mspage{389}

question \unsure{should|must} correspond to elliptic curves over $\mathbb Q$, equipped with a Jacobi rigidification of level $2$, \unsure{corresponding} (under the isomorphism $\deleted{\unsure{\mathcal T}}\unsure{\mathcal M}_{1,1\,\mathbb Q}'\simeq U_{0,3\,\mathbb Q}$) to the point $1/2$ of $U_{0,3}(\mathbb Q)$. The indeterminacy in the choice of such an elliptic curve is precisely in $H^1(\mathbb Q,\unsure{\{\pm1\}})$ (it is a \unsure{datum} of a quadratic extension \unsure{of|to} $\mathbb Q$) which is identified with $\mathrm{Hom}(\Gamma_{\mathbb Q},\deleted{\unsure{\mu}}\{\pm1\})$ -- so this is indeed also the indeterminacy for a lifting. This proves \unsure{essentially} that the idea suggested by the general yoga is indeed \unsure{correct|correct} -- and this shows at the same time that there is no (even on $\Gamma_{\mathbb Q,\sigma}$) \unsure{really} privileged choice. That there are such liftings \inserted{(on $\Gamma_{\mathbb Q,\sigma}$)} comes from the fact that there are indeed elliptic

\scanpage{181}

curves answering the question, the simplest being given by the equation
\[
(22)\qquad y^{\unsure{2}}-x(x-1)(x-1/2)=0
\]
-- \unsure{and} there is the choice made ! It would remain to try to make it explicit \deleted{by} in a computational \unsure{formulary}, if this is not \unsure{hopeless}. I should examine this \deleted{question} question more closely, sometime.

\scanpage{182}
\mspage{390}

\uline{46}\quad \uline{Definitive notations} (?) -- \uline{new recapitulation}

(\uline{for} the loop automorphisms of $\widehat{\Pi}_{0,3}$ and of $\widehat{\mathcal G}_{0,3}^{+}\simeq\unsure{\mathrm{Sl}(2,\mathbb Z)^{\prime\wedge}}$).

\textcircled{I}\quad \uline{The groups} $\widetilde M_{0,3}$ \uline{and} $\widetilde N_{0,3}^{!}$.

We set

$(1)\qquad\widetilde M_{0,3}=$ group of automorphisms of $\widehat{\mathcal G}_{0,3}^{+}\simeq\unsure{\mathrm{Sl}(2,\mathbb Z)^{\wedge\prime}}$ which leave $\widehat{\Pi}_{0,3}$ invariant, and which respect the ``loop structure'' i.e. \unsure{satisfy} $u(\unsure{\varepsilon_i|\varepsilon_0})=\operatorname{int}(f^{-1})(\varepsilon_i^\mu)$, with $f\in\widehat{\mathcal G}_{0,3}^{+}$, $i\in\{0,1,\infty\}$, $\mu\in\widehat{\mathbb Z}^{*}$ (OBSERVATION $f\in\widehat{\Pi}_{0,3}$, and $\unsure{f|f_0}$ is then determined in $L_0\backslash\widehat{\Pi}_{0,3}$)

$\overset{\mathrm{restr.}}{\overset{\sim}{\longrightarrow}}$ group of \deleted{outer} loop automorphisms of $\widehat{\Pi}_{0,3}$ which commute as outer actions with the natural outer action of $\unsure{\mathfrak S_3|\mathcal G_3}$.

\margin{\uline{N.B.} $\widetilde M_{0,3}$ is also the normalizer of $\widehat{\mathcal G}_{0,3}^{+}$ in $\widehat{\widehat{\mathcal G}}_{0,3}$, and likewise $\widetilde N_{0,3}$ the \unsure{normalizer|normalizer} of $\unsure{\widehat{\mathcal T}_{0,3}}=\unsure{\mathcal T_{0,3}}$ in $\widehat{\widehat{\mathcal T}}_{0,3}$}

We have an extension structure
\[
(2)\qquad 1\longrightarrow\widehat{\Pi}_{0,3}\longrightarrow\widetilde M_{0,3}\longrightarrow\widetilde N_{0,3}\longrightarrow1,
\]
\unsure{where} \deleted{$\unsure{N_0}$}
\[
(3)\qquad\widetilde N_{0,3}\overset{\mathrm{dfn}}{=}\widetilde M_{0,3}/\widehat{\Pi}_{0,3}
\]
$\simeq$ group of \uline{outer} loop automorphisms of $\widehat{\Pi}_{0,3}$, commuting with $\unsure{\mathfrak S_3|\mathcal G_3}$.

The natural homomorphism: $\unsure{\mathfrak S_3|\mathcal G_3}\longrightarrow\operatorname{Aut}(\unsure{\mathfrak S_3|\mathcal G_3})$ is an \unsure{isomorphism}
\[
(4)\qquad\unsure{\mathfrak S_3|\mathcal G_3}\overset{x\mapsto\operatorname{int}(x)}{\overset{\sim}{\longrightarrow}}\operatorname{Aut}(\unsure{\mathfrak S_3|\mathcal G_3}),
\]
and the natural homomorphism
\[
\widetilde M_{0,3}\longrightarrow\operatorname{Aut}(\unsure{\mathfrak S_3|\mathcal G_3}=\widehat{\mathcal G}_{0,3}^{+}/\widehat{\Pi}_{0,3})
\]
then defines a homomorphism
\[
(5)\qquad\widetilde M_{0,3}\longrightarrow\unsure{\mathfrak S_3|\mathcal G_3}
\]
which can also be interpreted as the composite
\[
(6)\qquad\widetilde M_{0,3}\hookrightarrow
\underset{\substack{\text{all the loop automorphisms}\\\text{of }\widehat{\Pi}_{0,3}}}{\widehat{\widehat{\mathcal G}}_{0,3}}
\xrightarrow{\substack{\text{action on the}\\\text{loop classes}}}
\unsure{\mathfrak S_3|\mathcal G_3},
\]